\section{Spatial regularity}\label{sec:regularidad}

We now fix the mass and move from one line to another in the same family. Translating a diagonal also changes the segment inside the
square. The estimate must therefore control both the variation of the factors and the effect of moving the endpoints of the segment.

We obtain regularity by comparing the profile with its translates using coercivity of a dual formulation. The argument requires no
derivatives of the maximizer. On $\cH=\bigoplus_fL^2(I_f)$ we use the norm $\|w\|_\cH^2=\sum_f\|w_f\|_{L^2(I_f)}^2$.

\subsection{Fixed activity}

Fix $\rho\in(0,1)$, $\lambda=\lambda(\rho)$, $\nu=\nu_\rho$ and the continuous representatives from Lemma~\ref{lem:vacancias}. We have
$u_f=1/(1+\lambda J_f(u))$, where $J_f(u)$ is the integral over the fiber of $\prod_{g\ne f}u_g\circ\ell_g$, $0\le J_f\le1$, and
$\kappa_\lambda:=(1+\lambda)^{-1}\le u_f\le1$. For bounded fields $v_f=\log u_f$, define \begin{equation}\label{eq:dual}
 \cE_\lambda(v)=\sum_f\int_{I_f}(e^{v_f}-v_f)+\lambda\int_{[0,1]^2}\exp\Bigl(\sum_fv_f\circ\ell_f\Bigr).
\end{equation} Its gradient in $\cH$ is $\cG_{\lambda,f}(v)=e^{v_f}-1+T_f\nu_v$, with $\nu_v=\lambda\exp(\sum_fv_f\circ\ell_f)$, and
the field of the maximizer satisfies $\cG_\lambda(v)=0$. The comparison fields are constrained by activity, with their mass free to
vary. Differentiation is justified by boundedness and dominated convergence along segments of bounded fields. For a bounded direction
$w$, \begin{equation}\label{eq:hess}
 \langle w,D\cG_\lambda(v)w\rangle=\sum_f\int e^{v_f}w_f^2+\int\nu_v\Bigl(\sum_fw_f\circ\ell_f\Bigr)^2 .
\end{equation}

\begin{lema}[Stability]\label{lem:estabilidad} Let $0<\kappa'\le1$. If $u=e^v$ and $\widetilde u=e^{\widetilde v}$ take values in
$[\kappa',1]$, then $\langle\cG_\lambda(\widetilde v)-\cG_\lambda(v),\widetilde v-v\rangle\ge\kappa'\|\widetilde v-v\|_\cH^2$. In
particular, if $\cG_\lambda(v)=0$, then $\|\widetilde v-v\|_\cH\le\|\cG_\lambda(\widetilde v)\|_\cH/\kappa'$. \end{lema}

\begin{proof} The exponentials along the segment between the logarithms lie in $[\kappa',1]$. Integrate~\eqref{eq:hess} along this
segment and apply Cauchy--Schwarz. \end{proof}

Extend $u_D$ and $u_A$ by one outside $I_D$ and $I_A$ (these extensions are continuous by Lemma~\ref{lem:vacancias}), and write
$\widehat u_D,\widehat u_A$. For $t\in\mathbb R$ let \[
 \widehat u_R(t)=\Bigl[1+\lambda\int_0^1u_C(c)\,\widehat u_D(t-c)\,\widehat u_A(t+c)\,dc\Bigr]^{-1},
\] which agrees with $u_R$ on $[0,1]$ and takes values in $[\kappa_\lambda,1]$. For a shift $t$, the comparison field on the original domains is \[
 u_R^t(r)=\widehat u_R(r+t),\quad u_C^t=u_C,\quad u_D^t(d)=\widehat u_D(d+t),\quad
 u_A^t(s)=\widehat u_A(s+t),\qquad v^t=\log u^t .
\]

\begin{lema}[Translation residual]\label{lem:residuo} With $\eta^t=\cG_\lambda(v^t)$, we have $\eta^t_R=0$,
$|\eta^t_C|,|\eta^t_D|,|\eta^t_A|\le2\lambda|t|$, and $\|\eta^t\|_\cH\le2\sqrt5\,\lambda|t|$. Consequently,
$\|v^t-v\|_\cH\le2\sqrt5\,\lambda(1+\lambda)|t|$. \end{lema}

\begin{proof} The first equality is the definition of $\widehat u_R$. In a column $c$, substitute $r'=r+t$ in the integral of
$u_R^t(r)\,u_D^t(r-c)\,u_A^t(r+c)$. The integrands agree with the original ones where the two domains overlap, and $r'\in[0,1]$ is
replaced by $r'\in[t,1+t]$, with the symmetric difference having length at most $2|t|$. The products of vacancies have absolute value
at most one. On the diagonal $d$, the same substitution gives $c=r'-(d+t)$ and the integrand is $\widehat u_R(r')\,u_C(c)\,\widehat
u_A(2r'-(d+t))$. The translated domain requires $r'\in[t,1+t]$ and $c\in[0,1]$, and the original domain of the diagonal $d+t$ requires
$r'\in[0,1]$ and $c\in[0,1]$, so their difference lies in the same two strips. If $d+t\notin I_D$, the original fiber is empty and
$\widehat u_D(d+t)=1$, and the argument still applies. For $A$ use $c=s+t-r'$. The residual is $\lambda u^t_f$ times the difference of
the two integrals, and its absolute value is at most $2\lambda|t|$. The domains of $C$, $D$, and $A$ have lengths $1$, $2$, and $2$,
respectively, giving the bound in $\cH$. The last inequality follows from Lemma~\ref{lem:estabilidad} with $\kappa'=\kappa_\lambda$.
\end{proof}

For translations in the second coordinate, interchange $R$ and $C$. Define $\widehat u_C$ by its equation and compare $u_R(r)$,
$\widehat u_C(c+t)$, $\widehat u_D(d-t)$, and $\widehat u_A(s+t)$. Component $C$ of the residual vanishes, and the other three satisfy
the same bounds.

To pass from $L^2$ to pointwise differences, suppose that $u=e^v$ and $\widetilde u=e^{\widetilde v}$ take values in $[0,1]$.
Telescoping the product of three factors on a fiber gives \begin{equation}\label{eq:conversion}
 |J_f(\widetilde u)-J_f(u)|\le\sum_{g\ne f}\|\widetilde u_g-u_g\|_{L^2(I_g)}\le\sqrt3\,\|\widetilde v-v\|_\cH ,
\end{equation} because on a fiber of $f$ the parameter $\ell_g$, $g\ne f$, varies with slope of absolute value one or two (never zero).
Cauchy--Schwarz and the parameter length $\le1$ give the first inequality, and $|e^\alpha-e^\beta|\le|\alpha-\beta|$ for
$\alpha,\beta\le0$ gives the second. For the translated field, $u^t_f(1+\lambda J_f(u^t))=1+\eta^t_f$. Comparing with $u_f(1+\lambda
J_f(u))=1$, using the bounds $\ge1$ on the denominators, Lemma~\ref{lem:residuo}, and~\eqref{eq:conversion}, we obtain
\begin{equation}\label{eq:traslacion}
 |u^t_f-u_f|\le2\lambda|t|+2\sqrt{15}\,\lambda^2(1+\lambda)|t|\le L(\lambda)|t|,\qquad
 L(\lambda)=2\lambda+8\lambda^2(1+\lambda).
\end{equation} Translation in the row coordinate controls $R$, $D$, and $A$, and translation in the column coordinate controls $C$.
